\section{Introduction}

Layer 0 suppressors sit at the first bottleneck in transformer models and regulate how much the network hedges when it is uncertain. Paper 1 showed that a small coalition of early heads systematically downweight factual continuations and boost hedging or editorial tokens \cite{thompson2025suppressors}. Ablating these heads improves factual preference and calibration which ties suppressor strength directly to hallucination behaviour. Paper 2 showed that structured noise is developmental fuel for circuit formation but only in sub saturated regimes and identified observable dependent saturation boundaries in the learning landscape \cite{thompson2025timescale}. In saturated regimes no amount of noise helps. Together these results described the anatomy and dynamics of a single phenomenon: how intelligence crystallizes from a noisy early phase under constraint.

Work on ironic rebound in transformer models lives one layer up the stack. Mann et al.\ show that when a model receives a negation instruction such as ``do not mention X'' it often becomes more likely to produce X than in a neutral baseline and that this effect strengthens with longer or more semantic distractors \cite{mann2025dontthink}. Their circuit analysis finds early heads that suppress forbidden tokens and sparse middle layer heads that amplify them under load. This reveals a dynamic negotiation between suppression and rebound in the middle of the network. In this paper I focus on the homeostatic story at layer 0. Instead of asking how models fail to obey negation under context load I ask how suppressor equilibria shape which circuits crystallize during grokking.

These papers left a central question open. Why is the system's preference for suppressors so strong. In parity grokking experiments I introduced an omega parameter that scales the output of a chosen layer 0 head. The natural setting $\omega = 1.0$ produced the slowest grokking. Perturbing omega in either direction to $0.5$ or $1.5$ accelerated grokking by roughly forty percent. The relationship between suppressor strength and grokking speed was U shaped rather than linear. If weakening suppressors helps and strengthening them also helps, why does the system settle on an intermediate value. Why does it resist learning most strongly at its natural configuration.

The shape of the curve suggests that the system is not merely following different paths through the same loss landscape. It points instead to an equilibrium. At $\omega = 1.0$ the network behaves as if it were in a metastable state. Training pushes on the weights but the system sits in a memorizing configuration for thousands of steps before collapsing into a generalizing circuit. Perturbing omega pushes the system away from this balance and induces instability. Instability accelerates escape from the memorization plateau.

This behaviour is reminiscent of Le Chatelier's principle. When a physical system at equilibrium is perturbed it responds by shifting internal variables in a way that opposes the perturbation. Heat a gas and it expands to reduce the change in temperature. Compress a gas and it develops pressure. The system works to restore balance \cite{kringelbach2024thermodynamics}. If neural networks obey an analogous principle then when I weaken a suppressor head other heads should adjust their outputs to preserve some quantity such as total information flow or attention budget. If I block that adjustment by freezing the other heads in place learning should be severely impaired.

The central hypothesis of this paper is that transformer networks maintain a homeostatic equilibrium through active variance redistribution across attention heads. When I perturb a suppressor by scaling its output the remaining heads adjust to oppose the change. This redistribution is not a passive by product of optimization. It is a coordinated response that preserves a preferred equilibrium state and shapes how the model traverses weight space during training.

To test this hypothesis I designed a kill test on a small transformer trained to grok the parity function. In the baseline condition $\omega = 1.0$ and all heads are free. In the perturbed condition $\omega = 0.5$ and all heads are free so other heads can compensate. In the frozen condition $\omega = 0.5$ and heads 1 through 3 in layer 0 are frozen at their initial values. Head 0 is perturbed and the remaining heads cannot compensate. If compensation is passive drift the frozen runs should grok at similar times to the perturbed runs and differ only in the path taken. If compensation is active and necessary the frozen runs should show severe impairment.

Across five seeds I find that blocking compensation always harms learning. In the cleanest seed frozen heads grok at $10{,}800$ steps compared to $2{,}200$ steps in the perturbed condition, a $4.9$ times slowdown. In other seeds the slowdown is smaller but consistently present. In one seed the frozen run groks fastest but converges to a brittle solution with test accuracy near $0.5$ while the free run groks more slowly but reaches $0.98$ test accuracy. This shows that compensation does not just control how fast the model learns. It shapes which circuits crystallize and whether they generalize beyond the training set.

These results bridge mechanistic interpretability and thermodynamic reasoning. Mechanistically I can identify which heads compensate when suppressors are perturbed and measure how freezing them alters grokking trajectories. Thermodynamically the pattern is consistent with a Le Chatelier style response: when I perturb suppressors the system shifts internal degrees of freedom to restore a preferred equilibrium. When I block those shifts the system either slows dramatically or learns the wrong thing. Understanding neural network learning therefore requires thinking in terms of homeostasis and equilibrium rather than only in terms of gradients and optima.
